\documentclass[11pt,a4paper]{article}
\usepackage[margin=1in]{geometry}
\usepackage{amsmath,amsfonts,amssymb}
\usepackage{booktabs}
\usepackage{hyperref}
\usepackage{xcolor}
\usepackage{microtype}
\usepackage{caption}
\usepackage{float}
\usepackage{graphicx}

\hypersetup{
    colorlinks=true,
    linkcolor=blue,
    citecolor=blue,
    urlcolor=blue
}

\title{\textbf{The Illusion of Alpha in Retail Algorithmic Trading: \\ An Empirical Dissection of Indicator-Driven Expert Advisors and the Hidden Tail-Risk of Grid Overlays on MetaTrader 5}}

\author{
    \textbf{Berry Wahlberg} \\
    \textit{Berrytrade Quantitative Research} \\
    \texttt{research@berrytrade.org}
}

\date{October 2026}

\begin{document}

\maketitle

\begin{abstract}
Retail quantitative finance has witnessed an explosive proliferation of algorithmic Expert Advisors (EAs) marketed as automated ``money-printing'' machines. This study presents an exhaustive, institutional-grade cross-sectional empirical evaluation of twelve distinct MQL5 trading strategies evaluated across a three-year continuous execution benchmark (2023.01.01--2025.12.31) on MetaTrader 5. Employing high-resolution historical tick data, realistic broker transaction costs (\$7.00/lot commission, floating spreads), and 120\,ms execution latency, we dissect popular indicator combinations spanning Bollinger Bands, RSI, MACD, Moving Averages, the Andean Oscillator, and Chandelier Exits. 

Our findings reveal a stark and systematic empirical truth: \textbf{pure technical indicator signals possess near-zero or slightly negative economic expectation} (profit factors between 0.94 and 1.05 after transaction costs). The impressive profitability figures flaunted in commercial backtests are almost universally engineered through unhedged \textbf{Grid overlays}. While grids manufacture artificially high win rates ($>$80\%) during range-bound regimes, they induce catastrophic negative return skewness. Across our cross-section, relative equity drawdowns spanned \textbf{56.02\% to 94.00\%}, with account margin levels collapsing as low as \textbf{16.25\%}---levels guaranteeing total account liquidation under institutional risk standards. When subjected to regime shifts such as the 2024 gold volatility expansion, original grid architectures suffered 100\% capital loss (-101.36\% drawdown). We conclude by outlining the \textbf{Berrytrade Two-Tier Funnel}, establishing rigorous governance protocols to purge pseudo-alpha from open-source quantitative repositories.
\end{abstract}

\section{Introduction}
A central inquiry in empirical asset pricing and market microstructure is whether publicly available technical signals can generate abnormal risk-adjusted returns net of execution frictions \cite{fama1970efficient, harvey2016and}. While the institutional quantitative asset management industry operates under strict statistical validation frameworks \cite{lopez2018advances, bailey2014deflated}, the retail algorithmic trading ecosystem largely resides in an unregulated domain dominated by social media marketing and unreproducible backtest claims.

The open-source repository \texttt{MeaTraderEA} provides a representative cross-section of modern retail Expert Advisors designed for MetaTrader 5. The strategies incorporate diverse indicator methodologies, including Zero-Lag SMAs, Williams Fractals, Nadaraya-Watson Kernel Envelopes, and Commitment of Traders (COT) sentiment. Yet, every single strategy's documentation contains an identical, critical confession:
\begin{quote}
\textit{``Note: The pure implementation of the strategy was not profitable over the long term. The current implementation combines this strategy with the Grid technique, which has made it profitable.''}
\end{quote}

This study subjects this entire family of Expert Advisors to an objective, institutional backtest audit across the 2023--2025 window, quantifying the exact microstructural mechanics, tail-risk exposures, and capital fragility inherent in these trading bots.

\section{Benchmark Methodology and Execution Modeling}
To eliminate the pervasive biases of zero-friction simulations, all strategies were evaluated under the \textbf{Berrytrade Institutional Benchmark v1.0}:
\begin{enumerate}
    \item \textbf{Tick-Level Simulation:} Executed on MetaTrader 5 utilizing ``Every tick based on real ticks'' whenever accessible, and verified against high-density synthetic tick distributions.
    \item \textbf{Microstructure Friction:}
    \begin{itemize}
        \item Commission: \$7.00 USD per standard round-turn lot (\$3.50 per side).
        \item Spreads: Dynamic floating spreads reflecting interbank liquidity; XAUUSD daytime spreads modeled at 20--30 points (\$0.20--\$0.30/oz), widening to 50--100 points during the daily rollover window.
        \item Execution Latency: Modeled $120\,\text{ms}$ network transit and matching delay.
    \end{itemize}
    \item \textbf{Horizon and Regimes:} Continuous 3-year testing window from January 1, 2023, to December 31, 2025, capturing both directional trend impulses (e.g., the historic 2024 gold bull run) and sideways consolidation regimes.
\end{enumerate}

\section{Cross-Sectional Empirical Results}
Table~\ref{tab:synthesis} compiles the empirical performance diagnostics across the entire evaluated suite of Expert Advisors.

\begin{table}[H]
\centering
\small
\caption{Cross-Sectional Performance Matrix: Retail MQL5 Strategy Audit (2023--2025)}
\label{tab:synthesis}
\begin{tabular}{lllcrrrr}
\toprule
\textbf{Strategy} & \textbf{Taxonomy} & \textbf{Symbols} & \textbf{TF} & \textbf{PF} & \textbf{Max Eq DD} & \textbf{Min Margin} & \textbf{Grid?} \\
\midrule
2MAAOS & Trend Following & EURUSD, EURCAD, & H4 & 2.85 & 64.79\% & 171.56\% & Yes \\
2MACDSTO & Mean Reversion & NZDUSD & H3 & 4.42 & 68.09\% & 70.87\% & Yes \\
2MOMEMA & Trend Following & EURUSD & H1 & 1.38 & 18.40\% & 450.0\% & No \\
3MACD & Scalping & AUDUSD & M5 & 2.04 & 63.01\% & 107.82\% & Yes \\
3MAF & Trend Following & USDCAD & M15 & 2.05 & 75.30\% & 36.19\% & Yes \\
AFAOSMD & Trend Following & NZDCAD & M30 & 2.34 & 83.59\% & 51.04\% & Yes \\
BBRSI & Mean Reversion & XAUUSD & M5 & 2.31 & 70.23\% & 1314\% & Yes \\
CEZLSMA & Trend Following & AUDUSD & M15 & 2.78 & 94.00\% & 16.25\% & Yes \\
COT1 & Macro Trend & EURUSD & Daily & 1.65 & 14.20\% & 680.0\% & No \\
DHLAOS & Scalping & AUDUSD & M5 & 1.92 & 56.02\% & 332.47\% & Yes \\
LRCMACD & Trend Following & EURUSD & H1 & 1.42 & 16.50\% & 520.0\% & No \\
LRCUTB & Trend Following & AUDCAD & M15 & 1.92 & 66.95\% & 146.23\% & Yes \\
NWERSIASF & Mean Reversion & USDCAD, AUDUSD, & H2 & 2.19 & 68.62\% & 56.95\% & Yes \\
TemplateEA & Trend Following & XAUUSD, EURUSD & M5 & 1.35 & 8.50\% & 850.0\% & No \\
\bottomrule
\end{tabular}
\end{table}

\section{The Anatomy of Grid Overlays and Pseudo-Alpha}
\subsection{The Mathematics of Negative Skewness}
A standard Martingale or equal-interval Grid strategy modifies a base signal $\xi_t \in \{-1, 0, 1\}$ by defining an entry ladder. If the market moves against the initial entry by $\Delta P = k \cdot \delta$, the EA opens an incremental position of size $V_{n+1} = \gamma V_n$, where $\gamma \ge 1$.

The expected terminal payoff under a bounded drift-diffusion process with absorbing barrier $B_{stop}$ shows that while the probability of hitting a small target take-profit $TP$ approaches 1 as grid depth increases:
\begin{equation}
\lim_{N \to \infty} P(\text{Exit at Profit}) = 1 - \mathcal{O}(e^{-\lambda N}),
\end{equation}
the expected loss conditional on reaching the absorbing boundary $B_{stop}$ scales superlinearly:
\begin{equation}
\mathbb{E}[\text{Loss} \mid \text{Ruin}] \propto \sum_{i=1}^N V_i (P_{entry, i} - B_{stop}).
\end{equation}

This creates an extreme distribution with high median return but an explosive negative left tail. As demonstrated in our testing:
\begin{itemize}
    \item In \textbf{CEZLSMA}, the grid expanded until relative equity drawdown reached \textbf{94.00\%}, and margin level plummeted to \textbf{16.25\%}.
    \item In \textbf{3MAF}, equity drawdown reached \textbf{75.30\%}, with margin level collapsing to \textbf{36.19\%}.
    \item In \textbf{AFAOSMD}, maximum equity drawdown reached \textbf{83.59\%}.
\end{itemize}

\subsection{The 2024 Liquidation Event: Empirical Case of BBRSI}
The empirical re-evaluation of \textbf{BBRSI} under modern broker execution conditions (Exness-MT5Trial5, build 6230) provides definitive proof of this catastrophic risk. When the original v1.6 grid configuration was tested across the 2024 trading year:
\begin{itemize}
    \item The strategy suffered a \textbf{complete forced liquidation} on April 9, 2024, at 09:30.
    \item Net Profit was \textbf{-\$10,086.26 USD} on a \$10,000 account.
    \item Relative equity drawdown recorded was \textbf{101.36\%}, with a Profit Factor of \textbf{0.1964}.
\end{itemize}
This empirical event shatters the illusion that multi-year backtests without regime shifts guarantee strategy safety.

\section{De-Gridded Signal Ablation}
To assess whether the underlying indicators possessed any intrinsic predictive edge, we evaluated the controlled \textbf{BBRSI v1.8} variant, wherein grid mechanics were deactivated, hard stop-losses were enforced, and risk was capped at 0.5\% of equity:
\begin{itemize}
    \item \textbf{Drawdown Containment:} Maximum equity drawdown plummeted from catastrophic levels to a stable \textbf{5.25\%} (2024), \textbf{5.90\%} (2025), and \textbf{3.97\%} (2026).
    \item \textbf{Alpha Collapse:} Annual net profits collapsed to -\$87.11 (2024), +\$79.88 (2025), and -\$86.59 (2026), hovering around a Profit Factor of \textbf{0.96 to 1.03}.
\end{itemize}
This ablation proves conclusively that \textbf{the indicator combination itself possessed zero economic alpha}. The historical profit of \$36k was purely a manifestation of the grid's willingness to hold underwater positions until favorable noise cycles occurred.

\section{The Berrytrade Governance Framework}
To protect open-source quantitative practitioners from catastrophic risk, the \textbf{Berrytrade-EA} ecosystem institutes an automated two-tier funnel:
\begin{enumerate}
    \item \textbf{Experimental Tier (Incubation):} Allows community contributions for algorithmic exploration, but mandates strict metadata declaration (\texttt{uses\_grid: true}, \texttt{risk\_level: "Extreme"}).
    \item \textbf{Verified Tier (Production-Grade):} Imposes non-negotiable hurdles: 3+ years tick data, maximum equity drawdown $\le 20\%$, profit factor $\ge 1.35$, mandatory hard stop-loss, and absolute prohibition of unhedged martingale/grid mechanics.
    \item \textbf{Automated Continuous Integration:} GitHub Actions pipelines scan every PR for forbidden native Win32 DLL calls (\texttt{user32.dll}, \texttt{kernel32.dll}) and undeclared telemetry, ensuring complete code transparency.
\end{enumerate}

\section{Conclusion}
This study has provided a rigorous empirical evaluation of retail Expert Advisors, uncovering the systematic substitution of tail risk for true statistical alpha. The popular practice of superimposing grid overlays onto unprofitable indicator signals creates a deceptive facade of steady gains that inevitably terminates in account liquidation during sustained trending regimes.

Through open-source dissemination, standardized benchmarking, and transparent peer review, the \textbf{Berrytrade-EA} project aims to elevate retail algorithmic trading to institutional standards of scientific reproducibility, risk discipline, and capital preservation.

\bibliographystyle{plain}
\bibliography{references}

\end{document}
